\documentclass[11pt]{article}
\usepackage[a4paper,margin=2.5cm]{geometry}

\setlength{\parindent}{0pt}
\setlength{\parskip}{0.6em}

\begin{document}

\begin{flushright}
August 28, 2026
\end{flushright}

\vspace{0.8cm}

Dear Editor,

We are pleased to submit our manuscript entitled ``Adaptive Loss Balancing for Physics-Informed Neural Networks in Hydrodynamic Inverse Problems: Design Principles from Full-Convergence Experiments'' for consideration for publication in \textit{Computer Methods in Applied Mechanics and Engineering} (CMAME).

Physics-informed neural networks (PINNs) have emerged as a promising framework for solving partial differential equations, but their training is hindered by multi-task gradient imbalance and causality violation. In this work, we systematically investigate how gradient-based loss balancing (ReLoBRaLo) and causal time-domain weighting interact with parameter priors in PINN inverse problems, using full-convergence (30,000-epoch) experiments on the 2D shallow water equations with a localized drainage sink, plus generalization studies on the Burgers equation and the Navier--Stokes cylinder wake.

Our main contributions are:

\begin{enumerate}
\item A design principle that adaptive loss balancing must exclude parameter priors to avoid re-ill-posing inverse problems, backed by a gradient-decoupling argument grounded in neural-tangent-kernel (NTK) theory.
\item Full-convergence experiments showing that balancing the prior drives the drainage-coefficient error from 0.047\% to 37.2\% (a 791-fold degradation), while fixed-weight priors preserve sub-percent parameter recovery.
\item A conditional view of causal training, which fails to transfer to Adam-only optimization across three PDE benchmarks, including strongly time-dependent vortex shedding.
\item A physically consistent finite-volume reference solution (Lax--Friedrichs) for the shallow-water equations, verified for mass conservation to within 0.1\% of the inflow.
\end{enumerate}

We believe this work is well suited to CMAME because it provides concrete, reproducible design principles for PINN-based hydrodynamic inverse problems, a topic of growing interest to the computational mechanics community. The reference solution is computed with a verified finite-volume solver, and the manuscript reports full-convergence and multi-seed statistics, in line with the numerical rigor expected by CMAME.

The manuscript is original, has not been published previously, and is not under consideration for publication elsewhere. All authors have read and approved the submission.

The source code is publicly available at \texttt{https://github.com/YuhaoPeng1103/ca-pinn}.

Thank you for considering our manuscript. We look forward to your response.

\vspace{0.8cm}

Sincerely,

\vspace{0.4cm}

Yuhao Peng\\
School of Mathematics and Physics, China University of Geosciences (Wuhan)\\
Wuhan, China\\
Email: \texttt{3136496634@qq.com}

\end{document}
